\cotefile{157-2}
\scanpage{163}\mspage{108}
\uline{Théorème-Scholie 2} Soit $\unsure{M|\mathcal M}$ une catégorie, munie de deux ensembles de flèches
\[
W,\mathrm{Fib}\subset\mathrm{Fl}(M).
\]

\deleted{On supposera que}\inserted{On fait les hypothèses suivantes} 1°) $W$ contient les \unsure{Iso|Isos}, \deleted{\unsure{et que}} si $f,g$ sont composables, si deux parmi $f,g,gf$ sont dans $W$, la troisième aussi.

2°) \deleted{\unsure{$\mathrm{Fib}$}} $M$ est stable par $\underset{\leftarrow}{\lim}$ finies, $\mathrm{Fib}$ \deleted{\unsure{est}} contient les \unsure{Isomorphismes|isomorphismes}, est stable par composition et par changement de base.
% slanted
\margin{Ces hypothèses \inserted{de stabilité} sur $\mathrm{Fib}$ ne serviront que dans II, et encore seulement (!) dans le cas A), repris in toto \unsure{sous} la forme A$'$ plus bas (p.~\unsure{112})}

\inserted{Pour $S$ dans $M$, $\boxed{\underline{\mathrm{Fib}}\,S}$ désigne la ss-catégorie pleine de $M_{/S}$ formée des $X\longrightarrow S$ qui sont des $\mathrm{Fib}$.}

On distingue deux \deleted{\unsure{conditions pour}} \inserted{versions différentes} d'un axiome \deleted{\unsure{(A)}} de factorisation d'une $f:X\longrightarrow S$ quelconque dans $M$ en $X\xrightarrow{i}\bar X\xrightarrow{\bar f}S$, avec $i\in W$, $\bar f\in\mathrm{Fib}$.

A) On se donne, pour \deleted{\unsure{tout}} $S$ dans $M$, un foncteur
\[
M_{/S}\longrightarrow M_{/S}\qquad X\longmapsto C(X)\quad(=\bar X),
\]
\deleted{avec $C(X)$ \unsure{et}} et un morphisme de foncteurs en $X$
\[
X\xrightarrow{i_X}C(X),
\]
avec
\[
i_X\in W,\qquad C(X)\in\mathrm{Ob}\,\underline{\mathrm{Fib}}\,S.
\]

B) Désignant par
\[
\mathrm{Cof}_0\subset\mathrm{Fl}\,M
\]
l'ens. des flèches de $M$ qui possèdent la propriété de relèvement à droite p.r. à $\mathrm{Fib}$, on suppose ``à la Quillen'' que toute \deleted{\unsure{application}} flèche $f:X\longrightarrow S$ se factorise en
\[
X\xrightarrow{i_{\deleted{\unsure{X}}}}C(\deleted{\unsure{\bar X}}X)\xrightarrow{\bar f}S,
\]
avec
\[
i\in\mathrm{Cofib}_0\cap W,\quad \bar f\in\mathrm{Fib}.
\]
% archive sheet number: 162
